\section{Veri Seti ve Ön İşleme}
\label{CH:Veri}

\subsection{AUTSL veri seti}

Çalışmada AUTSL veri seti kullanılmıştır~\cite{AUTSL2020}. Veri setinde 226 işaret,
43 işaretçi ve toplam 36.302 video bulunmaktadır. Videolar 512$\times$512 çözünürlükte
ve saniyede 30 karedir. Resmî bölmeler \texttt{train} (28.142 video, 31 işaretçi),
\texttt{val} (4.418 video, 6 işaretçi) ve \texttt{test} (3.742 video, 6 işaretçi)
şeklindedir. Bölmeler arasında ortak işaretçi bulunmadığı ve sınıf numaralarının üç
bölmede de aynı anlama geldiği ayrıca kontrol edilmiştir. Böylece ölçülen doğruluk,
modelin daha önce görmediği kişiler üzerindeki başarısını göstermektedir.

Kaggle üzerinden indirilen kopyada 11 videonun açılamadığı görülmüştür. Bu dosyalarda
MP4 kutusuna ait \texttt{moov} bilgisi eksiktir, yani dosyalar yarım inmiştir. İlgili
videolar atlanmış, kalan 28.132 eğitim, 4.418 doğrulama ve 3.741 test videosu ile
çalışılmıştır. Sınıf numarası ile Türkçe kelime eşlemesi
\texttt{SignList\_ClassId\_TR\_EN.csv} dosyasından okunmakta, uygulamada kelimeler
Türkçe yazımıyla (örneğin \sol cikmak\sag{} yerine \sol çıkmak\sag) gösterilmektedir.

\subsection{Nokta çıkarımı}

Her video karesi MediaPipe Tasks API'ye video modunda verilmiş, iki el için 21'er nokta
ve üst gövde için 7 nokta (burun, iki omuz, iki dirsek, iki bilek) alınmıştır. Tüm veri
setinin çıkarımı tek bir bilgisayarda yapılmış, sonuçlar her video için bir
\texttt{.npz} dosyası olarak saklanmıştır (toplam 1,1 GB). Ellerin bulunma oranı sol
elde yüzde 82, sağ elde yüzde 86 olarak ölçülmüş, gövde noktaları tüm karelerde
yakalanmıştır. Elin bulunmadığı karelerde ilgili koordinatlar sıfırlanmakta ve iki
boyutlu bir \sol el var mı\sag{} göstergesi özellik vektörüne eklenmektedir.

MediaPipe'ın sol/sağ el etiketi, gövde modelinin sol ve sağ bilek noktalarıyla
karşılaştırılmış ve 28.689 el-karede yüzde 98 oranında uyumlu bulunmuştur. AUTSL
videoları aynalanmamıştır; telefondaki analiz görüntüsü de aynalanmadan işlenmektedir.

\subsection{Normalizasyon ve sabit uzunluk}

Kişinin kameraya uzaklığı ve boyu modeli etkilemesin diye her karedeki noktalar iki
omzun orta noktasına göre ötelenmiş ve omuz genişliğine bölünmüştür:
\begin{equation}
  \hat{p}_i = \frac{p_i - \tfrac{1}{2}(s_{\text{sol}} + s_{\text{sağ}})}{\lVert s_{\text{sol}} - s_{\text{sağ}} \rVert}
\end{equation}
Burada $p_i$ bir noktanın görüntü koordinatı, $s$ ise omuz noktalarıdır. Omuzların
bulunamadığı karelerde bir önceki geçerli değer kullanılmaktadır. Telefon görüntüsü
kare olmadığı için Kotlin tarafında koordinatlar önce piksel birimine çevrilmekte,
böylece eğitimdeki kare görüntüyle ölçek farkı ortadan kalkmaktadır.

Normalizasyondan sonra her kare
$2\times21\times2 + 7\times2 + 2 = 100$ boyutlu bir vektörle temsil edilmektedir
(koordinatlar $x$ ve $y$; derinlik kullanılmamıştır). Videoların uzunluğu değişken
olduğu için (eğitim kümesinde ortalama yaklaşık 60 kare) her dizi doğrusal
interpolasyonla 64 kareye örneklenmektedir. Modelin girdisi bu nedenle
$64\times100$ boyutludur.

\subsection{Veri artırma}

Eğitimde dört tür artırma uygulanmaktadır. En önemlisi zamansal seyreltmedir: kareler
rastgele 1, 2 veya 3 adımda bir alınarak 30, 15 ve 10 FPS'lik kayıtlara benzetilir.
Bunun nedeni telefon kamerasının gerçek kullanımda saniyede yaklaşık 9--12 kare
işleyebilmesidir. Diğer artırmalar dizinin iki ucundan yüzde 10'a kadar kırpma,
$\pm10^\circ$ döndürme, yüzde 90--110 arası ölçek değişimi ve küçük bir kaydırmadır.
Artırma yalnızca eğitim kümesine uygulanmaktadır.
